% TrustFMI @ ACCV 2026 — "Trustworthy Foundation Models in Medical Imaging", workshop W01.
% Submission 2026-09-25. Full paper: 6-8 pages, references excluded.
%
% The style files are the official ACCV 2026 author kit, not plain LNCS: accv.sty
% raises a hard error if the class is not llncs, and refuses to load without `runningheads`.
%
% `review` mode is what the submission needs, and it does three things by itself:
%   * anonymises — it overrides \author, \titlerunning, \authorrunning and \institute, so the author
%     block is NOT hand-edited and AUTHORS.md stays out of the submission;
%   * turns on line numbers, which ACCV requires for review and removes at camera-ready;
%   * prints the paper ID on every page, which is mandatory — replace ***** once OpenReview issues it.
\documentclass[runningheads]{llncs}
\usepackage[review,year=2026,ID=*****]{accv}

\usepackage[T1]{fontenc}
\usepackage[utf8]{inputenc}
\usepackage{graphicx}
\usepackage{booktabs}
\usepackage{amsmath}
\usepackage{siunitx}
\usepackage{xcolor}
\usepackage{tikz}
\usetikzlibrary{arrows.meta}
\usepackage{microtype}
% pagebackref is recommended for the review version and explicitly not for camera-ready.
\usepackage[pagebackref,breaklinks,colorlinks,citecolor=blue]{hyperref}

% Draft-only marker: every result that is not final must remain conspicuous.
\newcommand{\provisional}[1]{%
  \begingroup\setlength{\fboxsep}{1pt}%
  \colorbox{yellow!35}{\textcolor{red!70!black}{#1}}%
  \endgroup}

\begin{document}

\title{The Text Beside the Image: Detection, Utility and Leakage
  for Trustworthy Multimodal Medical Data and Beyond}

% Anonymised automatically by accv.sty in `review` mode — do not fill these in for the submission.
\author{}
\institute{}

\maketitle

\begin{abstract}
Medical images are released with the reports that describe them, and protecting the image does
not protect the report. This paper measures the text component of such releases. We measure identifier detection, downstream utility and residual identity
leakage on the same documents, with the pseudonymisation policy as the variable under test: 15
detectors, three release conditions and four corpora of medical reports, legal judgments, news and
other genres, and e-mail, in German, English, Chinese and Arabic. A fixed 13-detector union
reaches a person sensitivity of \num{0.9998} at specificity \num{0.8686} on the medical reports,
\num{0.9958} at \num{0.8504} on the legal judgments, \num{0.9352} at \num{0.9318} on news and
other genres, and \num{0.9906} at \num{0.6235} on e-mail. With this ensemble, frequency matching
with a public name list recovers zero identities by alignment across the four corpora; the names it
got right were ones the detector missed, left in clear text. Cross-document linkage ranks the
correct person first for \SI{0.93}{\percent} of e-mail queries without training and
\SI{3.94}{\percent} with it, against $1/3697$ chance and \SI{71.98}{\percent} on unmodified text.
On the medical reports it recovers nothing without training and \SI{0.71}{\percent} of 138 queries
with it, against $1/207$ chance and a \SI{2.73}{\percent} ceiling on unmodified text.
\end{abstract}

\begin{keywords}
pseudonymisation \and clinical text \and text in multimodal releases \and privacy evaluation
\end{keywords}

\section{Introduction}\label{sec:introduction}

Clinical data are multimodal: a chest radiograph travels with its report, a cardiology study with
its discharge letter. Releasing either half requires both to be protected, since an image may need
defacing or removal of burned-in text while its report needs identifier replacement. Artificial
intelligence scales both tasks, and its errors vary by distribution, language, script and identifier
class. Packh\"auser et al.~\cite{packhauser2022reid} re-identified patients from de-identified chest
radiographs and priced a defence against diagnostic utility~\cite{packhauser2023anonymization};
pathological speech has been evaluated against privacy, utility and
fairness~\cite{arasteh2024speaker}; clinical text is instead kept with the model inside the
hospital~\cite{hou2025llama}. We ask what survives transformation of the text, and what protecting
it costs.

\emph{De-identification} reduces identifying information. Reversible \emph{pseudonymisation}
replaces an identifier with a realistic \emph{surrogate}; \emph{anonymisation} prevents protects individuals. 
Stable pseudonyms preserve longitudinal utility but hand an adversary a persistent handle.
European guidance regards keyed message authentication codes as
robust~\cite{enisa2021pseudonymisation}, and the health-informatics standard is revised
accordingly~\cite{din25237}.

Detection benchmarks ask whether an identifying span was
found~\cite{vats2026redact,jha2026piibench,uppala2026privacyfilter}; utility studies measure
downstream change~\cite{vakili2022utility,berg2020impact} or information
loss~\cite{trienes2024infolossqa}; re-identification studies measure attack success. Frameworks
combine technical performance, information loss and intrusion tests~\cite{mozes2021nointruder},
while surrogate systems and masking benchmarks cover single
pieces~\cite{eder2019emails,pilan2022tab,osborne2022bratsynthetic}. No study states together what one
policy buys and what it costs. Data-protection officers release documents and cases, not average
tokens, as one unchanged mention can identify someone after thousands were removed, so we measure
surviving findings, exposed cases and linkable entities rather than token rates alone.

Our design covers four corpora in German, English, Chinese and Arabic: clinical letters, legal
judgments, news and other genres, and e-mail. Every document receives the same detector comparison,
release conditions and measurements. Two predictions are under investigation: that detector precision
governs utility more strongly than the choice between surrogates and typed placeholders, and that
detection sensitivity moves the two attacks in opposite directions, because removing a name obstructs
context linkage while leaving a larger, more stable population for frequency matching. 

\section{Material and Methods}\label{sec:methods}

\subsection{Corpora}\label{sec:corpora}
% 0.6 p — table, plus the one caveat per corpus that changes how its results read.

\subsection{Conditions}\label{sec:conditions}
% 0.3 p — A/B/C; B and C differ in exactly one axis level.

\subsection{Detectors and Ensembles}\label{sec:detectors}
% 0.5 p — fifteen detectors, combination rules, sizes 1-3, and the token budget: script-dependent,
% and fixed rather than proportional for a reasoning model.

\subsection{Operating Points}\label{sec:operating-points}
% 0.4 p — fast against maximum on sensitivity and specificity; parallel cost; the 10 % band; why a
% sensitivity floor is not optional.

\subsection{Measurements}\label{sec:measurements}
% 0.45 p — detection, utility, leakage A2-A5 with threat models, stability. Exposure per document,
% per case and per entity.

\section{Results}\label{sec:results}

\subsection{Detection}\label{sec:results-detection}
% 0.9 p — cost/quality front; the fast cell is usually not the worse cell; nothing transfers;
% language is not a detail; consensus deletes classes.

\subsection{Utility}\label{sec:results-utility}
% 0.8 p — precision not policy; surrogates against placeholders; the multilingual caveat, which
% belongs here rather than in the discussion.

\subsection{Leakage}\label{sec:results-leakage}
% 0.9 p — leakage at each recall level with the denominator held constant; the attacks disagree;
% A4 and why Rank-1 misleads; what actually survives, per document, per case, per entity.

\subsection{Stability}\label{sec:results-stability}
% 0.4 p — first collision, fragmentation and drift on text; N2 failing in both directions at once.

\section{Discussion}\label{sec:discussion}

The strongest result is practical. A cheap ensemble is already a credible deployment on German
clinical letters. The fast operating point detects \SI{89.5}{\percent} of the identifiers we placed
in \SI{0.30}{s} per document, is \emph{more} specific than the maximum---0.977 against
0.967---and is 191 times cheaper, which buys the maximum 9.7 percentage points of sensitivity.
Under a precise detector, realistic pseudonymisation also keeps clinical information extraction
statistically indistinguishable from the original text. This is the combination a hospital needs:
measurable protection without destroying the task for which the data are shared.

Yet, the two error rates do not settle the choice. At the maximum point \SI{73.2}{\percent} of
letters leave with no surviving finding at all; at the fast point only \SI{6.5}{\percent} do, and
the median letter retains three. Note that this is the same pair of systems that differ by only
0.097 in sensitivity. Document-level exposure separates them where sensitivity and specificity
barely do, which is precisely the argument for reporting the denominator a release is actually
judged on.

The strongest language-model attacker falls from near-perfect ranking to chance in a ten-candidate
closed world. Yet, frequency matching and context linkage move in opposite directions with recall:
stable surrogates remove names but make detected identities easier to count and follow.

The large ensemble is promising but not final. On CARDIO:DE and TAB it appears strong on both error
rates, which no small ensemble achieves. Its Table~\ref{tab:operating-points} entries are explicitly
\provisional{provisional}; no quantitative claim relies on them.

For practice, use the fast point unless a release requires more. Audit consensus by identifier
class, split results by language, and report per patient. Rank-5 should accompany Rank-1, and
specificity should accompany sensitivity. Note that this requires better denominators, not a new
model.

Protecting images does not repair a report, or vice versa. Text failures are structured by language,
class, and document length. A mixed-language recall averages results differing by more than 40
percentage points and is unsuitable as a release criterion.

Limitations narrow the claim. We repeated multilingual runs after an inadequate non-Latin token
budget, but detection scoring retains truncated replies that the conditions drop. Surrogates depend
on detected spans, and CARDIO:DE's constructed gold lacks natural variation. Linkage runs on two
corpora only. Enron covers 263 promising configurations, selected elsewhere, rather than all 1,743.
No corpus marks public figures for stratified candidate ranking. These limits define where
replication is needed.

\section{Conclusion}\label{sec:conclusion}

Safe sharing of clinical data needs both modalities protected: a guarded image does not compensate
for a report that names the patient. This study measures the text half; the image half is a separate
literature.
Artificial intelligence makes both possible at scale, but its failures must be measured where data
are released. For text we priced detection, utility and re-identification on the same documents. A
large union ensemble replaces almost every person mention and holds both attacks near their chance
rates. Detector precision governs
utility more than the choice between surrogates and placeholders, and sensitivity understates
exposure per document, case and language. There is therefore no single anonymisation score: a
release should state its operating point, per-language errors, surviving cases, attacker knowledge
and stability. Our operating points and measurements are released with acceptance of this paper.


% No acknowledgements in the review copy: ACCV requires them omitted for anonymity.

\bibliographystyle{splncs04}
\bibliography{references}

\end{document}
